%============================ 6. CONCLUSIONS ==================================
\isec{\faFlagCheckered}{Conclusions}
\label{sec:conclusions}

\begin{keybox}{\faCheckCircle\ \ Findings}
\begin{itemize}
  \item The RoboMaster detected the floor tag, drove straight to it from 1.20\,m and stopped before it in 10 of 10
        trials, using a 5-state EKF predicted with a unicycle model and corrected by wheel odometry and by the tag's
        range and bearing.
  \item The final distance estimate agreed with the tape to \textbf{0.28\,cm MAE}, with a \textbf{range of 0.96\,cm}
        ($-0.48$ to $+0.48$\,cm), no measurable bias and $\sigma=0.32$\,cm, about the resolution of the tape itself.
  \item Wheel odometry alone ended 1.98\,cm off on average and up to 5.8\,cm, because of slip on the rug and pose
        glitches in the driver.
  \item The camera \emph{range} is what fixes the along-track error. With bearing only the error rises to 1.98\,cm,
        close to having no camera at all (2.53\,cm), which matches the observability argument of Lecture 12A.
  \item Prediction provides the continuous, uncertainty-aware estimate that the controller and the outlier gate rely
        on, and it carries the estimate through camera gaps; correction keeps it from drifting.
\end{itemize}
\end{keybox}

\lead{What made the difference in practice.} Three problems were found and fixed before the trials, each visible in
the logs: the \code{/odom} twist field is late and too large (speed is now taken from the pose), the pose has rare
impossible jumps (now skipped by a plausibility guard), and a 154\,mm flat tag is unreliable beyond 1.2\,m with this
camera (start moved to 1.20\,m, stop to 0.50\,m).

\lead{Next steps.} Tighten $R_{tag}$ to about 1\,cm so that the reported $\sigma_d$ matches the observed spread; add a
final slow approach to remove the 1.5 to 2\,cm coasting at the stop; check absolute accuracy with an instrument independent
of the calibration (for example a second camera or a laser distance meter).

%============================ APPENDIX ========================================
\appendix
\isec{\faGithub}{Code, Data and Reproduction}
\label{sec:repro}

\begin{table}[H]
\centering\small
\rowcolors{2}{kfupmlight}{white}
\begin{tabularx}{\linewidth}{@{}l>{\raggedright\arraybackslash}X@{}}
\rowcolor{kfupm}\hd{Path} & \hd{Content}\\
\code{src/rm_odom_ekf/rm_odom_ekf/ekf.py} & the EKF: $f$, $F$, $G$, $H_{odom}$, $h$, $H_{tag}$, gate, update, wrap\\
\code{.../phase3_node.py} & ROS node: predict, odometry update, camera update, control, logging\\
\code{.../controller.py}, \code{tag_detector.py} & controller states; AprilTag range and bearing\\
\code{src/rm_odom_ekf/config/robot.yaml} & every parameter used in the trials\\
\code{src/rm_odom_ekf/test/} & 23 unit tests, including central-difference checks of $F$, $G$, $H_{tag}$\\
\code{results/trials/} & 10 trials: \code{steps.csv} (20\,Hz) and \code{summary.json}\\
\code{results/test_runs/}, \code{results/calibration/} & test drives, repeated first attempts, calibration data\\
\code{scripts/analyze_trials.py} & Table~\ref{tab:trials}, statistics, Figures~\ref{fig:scatter}--\ref{fig:stats}\\
\code{scripts/replay_ekf.py} & sensor-removal study, Table~\ref{tab:replay}\\
\code{theory/FORMULA_MAP.md} & every equation of this report and the code line that implements it\\
\code{video_tools/} & demo video: AR overlay, 3D replay, dashboard\\
\end{tabularx}
\caption{Repository: \link{\repoURL}{\repoURL}. Demo video: \link{\videolink}{\videolink}.}
\end{table}

\begin{lstlisting}[style=sh]
# unit tests (no robot needed)
cd src/rm_odom_ekf && python3 -m pytest -q test          # 23 passed
# one trial on the robot (driver running)
ros2 launch rm_odom_ekf phase3.launch.py trial:=N
python3 scripts/add_measurement.py N 0.492              # tape reading in metres
# statistics, figures and the replay study
python3 scripts/analyze_trials.py
python3 scripts/replay_ekf.py
\end{lstlisting}
